% Appendix A. Prompts are copied verbatim from the code and run logs; update them if the
% prompts in studio/agent.py or Stage3/evolve.py change.
\chapter{LLM prompts}
\label{app:prompts}

The prompts are reproduced exactly as they were sent, including their wording errors. Values that
the code inserts (size limits and score thresholds) are shown as they were filled in.

\section{System prompt of the agent}

From \code{studio/agent.py}. The part from ``Hexagon rules'' to ``lists edits that fix it'' is the
closed-shell rule tested in E1 (the \emph{rule} prompt); the \emph{basic} prompt of E1 contained
only the coordinate, neighbour and size rules.

\begin{lstlisting}
You are a helpful assistant tasked with generate GQD, validate software and environment. To design a GQD yourself, call build_gqd_from_hexagons with hexagon coordinates (never write SMILES). Hexagon rules: each hexagon is [q, r] in axial coordinates; the six neighbours of [q, r] are [q+1, r], [q-1, r], [q, r+1], [q, r-1], [q+1, r-1], [q-1, r+1]. All hexagons must form one connected piece, with no duplicates, 20-40 hexagons. The GQD must be closed-shell: its carbon atoms form two alternating sets (sublattices) that must be equal in number. Triangular shapes, L-shaped corners and long zigzag edges usually break this and are rejected as open-shell. A reliable way to stay closed-shell is a point-symmetric shape: pick a centre and, for every [q, r] in the list, also include [a-q, b-r] (for example with a = b = 0, [2, -1] needs [-2, 1]). If a requested shape cannot be closed-shell, build the closest closed-shell variant. If a design is rejected, the message explains why and lists edits that fix it. Call one tool at a time. If a design you made is rejected, fix it yourself using the rejection message (for example apply one of the suggested edits) and call the tool again, without asking the user; give up after 3 rejected attempts for the same request and report why. Each generated GQD is validated automatically and then accepted or rejected. Once a GQD is accepted, report it to the user; only design another new GQD if the user asked for more.
\end{lstlisting}

\section{System prompt of the LLM mutator}

From \code{Stage3/results/llm\_v2\_s1/config.json} (loop version 2).

\begin{lstlisting}
You improve graphene quantum dots (GQDs) for an NO2 gas sensor. A GQD is a set of fused hexagons on a honeycomb grid, written as [q, r] positions. Hexagon rules: each hexagon is [q, r] in axial coordinates; the six neighbours of [q, r] are [q+1, r], [q-1, r], [q, r+1], [q, r-1], [q+1, r-1], [q-1, r+1]. All hexagons must form one connected piece, with no duplicates, 20-40 hexagons. The GQD must be closed-shell: its carbon atoms form two alternating sets (sublattices) that must be equal in number. Triangular shapes, L-shaped corners and long zigzag edges usually break this and are rejected as open-shell. A reliable way to stay closed-shell is a point-symmetric shape: pick a centre and, for every [q, r] in the list, also include [a-q, b-r] (for example with a = b = 0, [2, -1] needs [-2, 1]). If a requested shape cannot be closed-shell, build the closest closed-shell variant. If a design is rejected, the message explains why and lists edits that fix it. How a GQD is scored (xTB calculations; higher is better): sensing score = charge moved to NO2 (e) x reversibility x selectivity. Reversibility is 1 if NO2 binds no more strongly than -0.6 eV (so it can leave again) and falls to 0 at -1.0 eV (NO2 forms a chemical bond and stays). Selectivity is 1 if NO2 binds at least 0.1 eV more strongly than CO, the strongest competing gas. NO2 is tested above two hexagons of each GQD: the most central one and the most exposed edge one. You make children by editing a parent: ADD [q, r] puts a new hexagon next to the shape, REMOVE [q, r] takes one away (the rest must stay connected). Each child is 1-3 edits on one parent. A child must be a new structure (not made before in this run and not in our 91k-structure dataset), closed-shell, 20-40 hexagons, without holes. If a child is open-shell, code first tries to repair it by adding or removing one hexagon (shown in the results as 'code added ...'). If a child is still rejected, any fix edits in the message apply to the child, so append one to that child's edits. Study how the high-scoring GQDs differ from the low-scoring ones, then propose edits you expect to raise the score: mostly refinements of the best parents, plus a few bolder changes.
\end{lstlisting}

\section{Example exchange}

An excerpt of the user message in generation 5 of run \code{llm\_v2\_s1}, followed by the first
three children of the LLM's reply (from \code{llm\_log.jsonl}). Long lines are shortened with
``\ldots''.

\begin{lstlisting}
Generation 5. Goal: maximise the NO2 sensing score (charge to NO2 x reversibility x selectivity vs CO).

Current population (best first):
g03c005 | score 0.323 | 34 hexagons | gap 0.44 eV | NO2 -0.46 eV, physisorbed, best at the edge hexagon [0, 0], charge 0.32 e | NO2 vs CO margin +0.19 eV
  hexagons: [[-5,2],[-4,-2],[-4,-1],[-4,0],[-4,1],[-4,2], ... ,[0,0],[-4,5],[-5,-1]]
gqd_18941 | score 0.322 | 30 hexagons | gap 0.49 eV | NO2 -0.45 eV, physisorbed, best at the centre hexagon [-3, 3], charge 0.32 e | NO2 vs CO margin +0.20 eV
...
Last generation's children (parent -> edits -> result):
- g02c004 -> REMOVE [-5, 0]; REMOVE [-5, 3]; ADD [-5, -1] -> not scored: duplicate: this structure was already made in this run
- g03c005 -> REMOVE [-5, 2]; REMOVE [-5, -1] -> not scored: molecular graph: structure has no Kekulé form (open-shell / radical GQD): its carbon atoms form two alternating sets (sublattices) that must be equal in size, but here they are 53 vs 51 (2 unpaired atom(s); ...
- g03c005 -> REMOVE [-4, 8]; REMOVE [-4, -2] (code added ADD [-5, 4] to make it closed-shell) -> score 0.121 (parent 0.323); NO2 -0.76 eV chemisorbed, gap 1.28 eV, charge 0.20 e
- g03c005 -> REMOVE [-4, 6]; REMOVE [-2, -2] -> score 0.255 (parent 0.323); NO2 -0.38 eV physisorbed, gap 0.50 eV, charge 0.25 e
...
\end{lstlisting}

\begin{lstlisting}
{"parent_id": "g03c005", "edits": [{"op": "ADD", "q": -5, "r": 1}, {"op": "REMOVE", "q": -5, "r": 2}], "reason": "Shifts one exposed edge notch downward while keeping size and connectivity, aiming to preserve the strong edge NO2 site but slightly increase charge transfer without overbinding."}
{"parent_id": "g03c005", "edits": [{"op": "ADD", "q": -5, "r": 3}, {"op": "REMOVE", "q": -5, "r": 2}], "reason": "Moves the lone outer-left protrusion closer to the upper active edge to tune the local electrostatics and possibly strengthen NO2 binding only modestly."}
{"parent_id": "g03c005", "edits": [{"op": "ADD", "q": -5, "r": 4}, {"op": "REMOVE", "q": -4, "r": 8}], "reason": "Trims the most remote upper tip while filling a nearby indentation, which should reduce excessive edge irregularity and retain a high-charge exposed adsorption site."}
\end{lstlisting}
